\documentclass[10pt,a4paper]{amsart}
\usepackage[margin=19mm]{geometry}
\usepackage{amsmath,amssymb,mathtools}
\usepackage[scr=rsfs,cal=euler]{mathalfa}
\usepackage{xfrac}
\usepackage[unicode,hidelinks,bookmarks=false]{hyperref}
\newcommand{\ie}{i.e.\ }
\newcommand{\cf}{cf.\ }
\newcommand{\resp}{respectively\ }
\newcommand{\Subsection}{\subsection}
\providecommand{\nobreakdash}{\nobreak\hbox{-}}
\setlength{\emergencystretch}{2em}
\multlinegap0pt
\usepackage{bm}
\usepackage{bbm}
\usepackage{dsfont}
\usepackage{xspace}
\usepackage{cancel}
\usepackage{mathtools}
\usepackage{amsthm}
\makeatletter
\@ifundefined{theorem}{\newtheorem{theorem}{Theorem}}{}
\@ifundefined{corollary}{\newtheorem{corollary}[theorem]{Corollary}}{}
\@ifundefined{definition}{\newtheorem{definition}[theorem]{Definition}}{}
\@ifundefined{lemma}{\newtheorem{lemma}[theorem]{Lemma}}{}
\@ifundefined{proposition}{\newtheorem{proposition}[theorem]{Proposition}}{}
\makeatother
\newcommand{\Bf}{\boldsymbol{f}}
\newcommand{\E}{\mathcal{E}}
\newcommand{\RR}{\mathbb{R}}
\newcommand{\SP}{\mathcal{P}}
\newcommand{\Supp}{\operatorname{Supp}}
\newcommand{\ZZ}{\mathbb{Z}}
\newcommand{\fa}{\mathfrak{a}}
\newcommand{\fpt}{{\operatorname{fpt}}}
\newcommand{\m}{\mathfrak{m}}
\newcommand{\tE}{\operatorname{\widehat{\E}}}
\title[$F$-pure thresholds and $F$-Volumes of some non principal id]{$F$-pure thresholds and $F$-Volumes of some non principal ideals}
\author{W\'agner Badilla-C\'espedes, Edwin Le\'on-Cardenal}
\date{Source: arXiv:2305.00571v3 (CC BY 4.0)}
\begin{document}
\maketitle

\noindent\emph{Source and licence.} $F$-pure thresholds and $F$-Volumes of some non principal ideals. Version arXiv:2305.00571v3 (CC BY 4.0), distributed under the Creative Commons Attribution 4.0 International licence, as recorded in the arXiv licence metadata for this exact version and shown on the abstract page https://arxiv.org/abs/2305.00571v3. Full rights notice: licenses/arxiv-2305.00571-CC-BY-4.0.txt.\par\smallskip

\noindent\textit{Editorial scope.} Counted unit: the Theorem whose source label is \texttt{Thm:F-pure-threshold-polytope} in the article (source0.tex lines 582--593, source label \texttt{Thm:F-pure-threshold-polytope}), delivered together with the article's own complete original proof of it (source0.tex lines 594--650, one \texttt{proof} environment of 57 lines). No mathematics is written, repaired or re-derived: statement and proof are the article's own text line for line. Required context retained uncounted: Cor:adding wo carrying (lines 265--269); Def:ExpMat\_SpliPol (lines 299--306); equal-vector (lines 372--379); propoBinomio (lines 394--400); Prop:Properties ftt ideals (lines 567--572). Every definition, display and label of those lines is retained unchanged. Omitted from this package and not redistributed: 1--264, 270--298, 307--371, 380--393, 401--566, 573--581, 651--1069. Nothing inside a retained excerpt is omitted or altered: every retained body is delivered exactly as the article's lines are written, so no omission receipt of any kind accompanies this package. Citations: the retained excerpts carry no citation command at all, so no bibliography entry of the article is transcribed and no reference is redistributed. Reference adaptation: none was needed and none was made; every reference inside the delivered unit resolves inside it, so no unresolved reference or citation is produced. Printed slips kept literal: no printed word, symbol, hypothesis or number of the counted statement or of its proof was altered. Layout and class: the article's own class is \texttt{amsart} with packages including amsmath, amsfonts, amssymb, setspace, moreverb, euscript, bbold, inputenc, fontenc, tikz, tikz-3dplot, wrapfig, pgfplots, color; the wrapper is the lane's pinned \texttt{amsart} editorial wrapper at 10pt on A4, which restates the theorem-like environments the retained text needs and adds the article's own macro definitions that the retained text needs (\texttt{\textbackslash Bf}, \texttt{\textbackslash E}, \texttt{\textbackslash RR}, \texttt{\textbackslash SP}, \texttt{\textbackslash Supp}, \texttt{\textbackslash ZZ}, \texttt{\textbackslash fa}, \texttt{\textbackslash fpt}, \texttt{\textbackslash m}, \texttt{\textbackslash tE}), with extra wrapper packages (bm, bbm, dsfont, xspace, cancel, mathtools). The delivered numbering is the unit's own. Assets: no retained span contains a figure environment, a graphics-inclusion command, a PDF-page inclusion, a table or a TikZ picture; no image file is copied into this package, the package declares no asset dependency and no image file is redistributed with it. Licence: version arXiv:2305.00571v3 of this article is distributed under the Creative Commons Attribution 4.0 International licence, as recorded in the arXiv licence metadata for this exact version and shown on the abstract page https://arxiv.org/abs/2305.00571v3. A full rights notice is delivered at licenses/arxiv-2305.00571-CC-BY-4.0.txt. Characters: every retained excerpt is pure ASCII, so the delivered engine, XeLaTeX with its default Latin Modern text font, reports no missing character.
\begin{corollary}\label{Cor:adding wo carrying}
	The entries of a vector $(\alpha_1,\ldots,\alpha_m)\in[0,1]^m$ add without carrying mod $p$ if and only if, for every $e\geq 1$
	\[\binom{p^e |\langle\alpha\rangle|}{p^e \langle\alpha\rangle_e}  \not \equiv 0 \bmod p.
	\]
\end{corollary}

\begin{definition}\label{Def:ExpMat_SpliPol}
Given a polynomial mapping $\Bf=(f_1, \ldots, f_t)$ we set $N=l_1+\cdots+l_t$ and define the following.
\begin{enumerate}
\item An \textit{exponent matrix} of $\Bf$ is an $(m \times N)$-matrix, lets say $\E_{\Bf}$, having as columns the different elements of $\bigcup_{i=1}^t\Supp(\hat{f}_i)$, i.e., $\E_{\Bf}=(\tE_{1}\mid \cdots\mid \tE_{t})$.
\item The \textit{splitting polytope} of $\Bf$ is the set 
\[\SP_{\Bf}=\{\gamma \in \RR_{\geq 0}^{N} \mid \E_{\Bf}\cdot \gamma \preceq \textbf{1}_{m}\}.\]
\end{enumerate}
\end{definition}

\begin{lemma}\label{equal-vector}
Assume that $\SP_{\Bf}$ has a unique maximal point $\rho\in [0,1]^N$, which is written as  $\rho=(\rho_1, \ldots, \rho_t)$ with $\rho_i \in \RR^{l_i}$, according to Definition \ref{Def:ExpMat_SpliPol}. Then the following statements hold.
\begin{enumerate}
\item If for some $\gamma \in \RR_{\geq 0}^N$, $|\gamma|=|\langle \rho \rangle_e|$ and $\E_{\Bf} \gamma = \E_{\Bf}  \langle \rho \rangle_e$,  then $ \gamma = \langle \rho \rangle_e $. Moreover, if 
there exist vectors $\phi, \psi \in \RR_{\geq 0}^{N}$, with $\psi \preceq \langle \rho \rangle_e$ and such that $|\phi|=|\psi|$ and  $\E_{\Bf} \phi = \E_{\Bf} \psi$;  then $\phi=\psi$.
\item Let $i$ be a fixed index in $\{1,\ldots,t\}$. If for some $\delta \in \RR_{\geq 0}^{l_i}$, $|\delta|=|\langle \rho_i \rangle_e|$ and $\tE_i \delta = \tE_i  \langle \rho_i \rangle_e$, then $ \delta = \langle \rho_i \rangle_e $. Moreover, if  there exist vectors $\phi, \psi  \in \RR_{\geq 0}^{l_i}$, with $\psi \preceq \langle \rho_i \rangle_e$ and such that $|\phi|=|\psi|$ and  $\tE_i \phi = \tE_i \psi$;  then $\phi=\psi$.
\end{enumerate}
\end{lemma}

\begin{proposition} \label{propoBinomio}
	 Under the assumptions of Lemma \ref{equal-vector}, the following statements hold.
\begin{enumerate}
\item The coefficient of the monomial $x^{p^{e}\tE_i \langle \rho_i \rangle_e}$ in $f_i^{p^{e} |\langle \rho_i \rangle_e|}$ is $\binom{p^{e}|\langle \rho_i \rangle_e|}{p^{e}\langle \rho_i \rangle_e}c_{a(i)}^{b_{e}(i)}$, where $b_{e}(i)$ is the vector in $\ZZ_{\geq 0}^{n_i}$ having $p^{e}\langle \rho_i \rangle_e$ in the first $l_i$ entries and $0$ in the remaining ones. Moreover, the coefficient of the monomial $x^{p^{e}\E_{\Bf} \langle \rho \rangle_e}$ in the polynomial $g:=\prod_{i=1}^{t} f_i^{p^{e}|\langle \rho_i \rangle_e|}$ is $\prod_{i=1}^{t}\binom{p^{e}|\langle \rho_i \rangle_e|}{p^{e}\langle \rho_i \rangle_e}c_{a(i)}^{b_{e}(i)}$.
\item Let $v$ be an element of  $\frac{1}{p^{e}}\ZZ_{\geq 0}^N$, satisfying $v \preceq \langle \rho \rangle_e$ and written in the form $v=(v_1, \ldots, v_t)$ with $v_i \in \frac{1}{p^{e}}\ZZ_{\geq 0}^{l_i}$. Then, the coefficient of the monomial $x^{p^{e}\tE_i v_i}$ in $f_i^{p^{e}|v_i|}$ is $\binom{p^{e}|v_i|}{p^{e} v_i}c_{a(i)}^{d_{e}(i)}$, where $d_{e}(i)$ is a vector of the form $p^e(v_i,0,\ldots,0)\in\ZZ_{\geq 0}^{n_i}$. Moreover, the coefficient of the monomial $x^{p^{e}\E_{\Bf} v}$ in the polynomial $g=\prod_{i=1}^{t} f_i^{p^{e}|v_i|}$ is $\prod_{i=1}^{t}\binom{p^{e}|v_i|}{p^{e}v_i }c_{a(i)}^{d_{e}(i)}$.
\end{enumerate}
\end{proposition}

\begin{proposition}\label{Prop:Properties ftt ideals}The following properties hold.
	\begin{enumerate}
		\item $\fpt(\fa)\leq \min\{t, \fpt(\fa^\circ)\}$.
		\item $\fpt(\fa^\circ)=\max_{\gamma\in \SP_{\Bf} } |\gamma|.$
	\end{enumerate}
\end{proposition}

\begin{theorem} \label{Thm:F-pure-threshold-polytope}
	Let $\fa \subseteq \m$ be a nonzero ideal with term ideal $\fa^\circ$, and take a list $\Bf=(f_1,\ldots,f_t)$  of minimal generators of $\fa$. Suppose that $\SP_{\Bf}$ has a unique maximal point $\rho=(\rho_1, \ldots, \rho_t)$ with $\rho_i \in \RR^{l_i}$.
	
	\noindent If $I$ denotes the set of indices for which $S_i$ is finite, then the following assertions hold.
	\begin{enumerate}
		\item If $I= \emptyset$ then $\fpt(\fa)=\fpt(\fa^\circ)$. 
		\item If $I \ne \emptyset$ then 
		\[
			\fpt(\fa) \geq \sum_{i \in \{1,\ldots,t\} \setminus I}|\rho_i| + \sum_{i \in I}\left(|\langle \rho_i \rangle_{S_i}|+\frac{1}{p^{S_i}}\right).
		\]
	\end{enumerate}
\end{theorem}

\begin{proof}
	For the case $I= \emptyset$, it will be enough to show that $|\rho| \leq \fpt(\fa)$, according to Proposition~\ref{Prop:Properties ftt ideals}. Our first goal is to find an exponent $Q(e)$ such that $\fa^{Q(e)} \nsubseteq  \m^{[p^{e}]}$. In order to check that this exponent is $p^{e} |\langle \rho \rangle_e|$ we start by recalling from Proposition~\ref{propoBinomio}(1) that
	\begin{equation}\label{Eq:Monomial in g}
		\prod_{i=1}^{t}\binom{p^{e}|\langle \rho_i \rangle_e|}{p^{e}\langle \rho_i \rangle_e}c_{a(i)}^{b_{e}(i)}x^{p^{e}\E_{\Bf} \langle \rho \rangle_e}
	\end{equation}
	appears as a summand of the polynomial  $g=\prod_{i=1}^{t} f_i^{p^{e}|\langle \rho_i \rangle_e|}$. Since $I=\emptyset$,  Corollary \ref{Cor:adding wo carrying} implies $\binom{p^{e}|\langle \rho_i \rangle_e|}{p^{e}\langle \rho_i \rangle_e} \not \equiv 0$ mod $p$ for each $\rho_i$, making the coefficient of \eqref{Eq:Monomial in g} not zero and thus $p^{e}\E_{\Bf} \langle \rho \rangle_e \in \Supp(g)$. However, from $\langle \rho \rangle_e \prec \rho$ we have 
	\[p^{e}\E_{\Bf} \langle \rho \rangle_e \prec p^{e}\E_{\Bf} \rho \preceq p^{e}\textbf{1}_m,
	\]
	meaning that  $x^{p^{e}\E_{\Bf} \langle \rho \rangle_e}$ does not belong to 
	 $\m^{[p^{e}]}$. In turn this implies that $g \notin \m^{[p^{e}]}$ and finally $\fa^{p^{e} |\langle \rho \rangle_e|} \nsubseteq  \m^{[p^{e}]}$ because $g$ was already a member of $\fa^{p^{e}|\langle \rho \rangle_e|}$. We conclude that $|\langle \rho \rangle_e| \leq \frac{\nu_{\fa}(p^{e})}{p^{e}}$, which take us to $|\rho| \leq \fpt(\fa)$ by taking the limit as $e \to \infty$.
	
	We now address the second part of the proof. We assume without lost of generality that $I=\{1,\ldots,r\}$ and define $S=\max_{1 \leq i \leq r}S_i$.  
	Our strategy here is to find again a suitable exponent $Q(e)$ verifying $\fa^{Q(e)} \nsubseteq  \m^{[p^{e}]}$, but in this case the construction of the corresponding monomial is more involved. Note first that the definition of $S_i$ gives
\begin{equation*}
\begin{matrix}
\rho_{1,1}^{(S_1+1)}+ \cdots + \rho_{1,l_1}^{(S_1+1)} \geq p\\
\rho_{2,1}^{(S_2+1)}+ \cdots + \rho_{2,l_2}^{(S_2+1)} \geq p\\
\vdots \\
\rho_{r,1}^{(S_r+1)}+ \cdots + \rho_{r,l_r}^{(S_r+1)} \geq p.
\end{matrix}
\end{equation*}
There should exist then non negative integers $n_{1,1}, \ldots, n_{1,l_1},$  $n_{2,1}, \ldots, n_{2,l_2},$ $\ldots, n_{r,1}, \ldots, n_{r,l_r}$ such that
\begin{equation*}
\begin{matrix}
n_{1,1}+ \cdots + n_{1,l_1} = p-1 &  \mathrm{and}& 0 \leq n_{1,j_1} \leq \rho_{1,j_1}^{(S_1+1)}\\
n_{2,1}+ \cdots + n_{2,l_2}= p-1 &  \mathrm{and}&0 \leq n_{2,j_2} \leq \rho_{2,j_2}^{(S_2+1)}\\
\vdots & \vdots & \vdots\\
n_{r,1}+ \cdots + n_{r,l_r} = p-1 &  \mathrm{and}& 0 \leq n_{r,j_r} \leq \rho_{r,j_r}^{(S_r+1)},
\end{matrix}
\end{equation*}
with at least one index in each row verifying the strict inequality, say $n_{i,1}<\rho_{i,1}^{(S_i+1)}$. 
For $e$ big enough we define the vectors, 
\[
\begin{cases}
		v_1=\langle \rho_1 \rangle_{S_1}+\left(\frac{n_{1,1}}{p^{S_1+1}}+\frac{p-1}{p^{S_1+2}}+ \cdots + \frac{p-1}{p^{S_1+e}},\frac{n_{1,2}}{p^{S_1+1}}, \ldots, \frac{n_{1,l_1}}{p^{S_1+1}} \right)\\
		v_2=\langle \rho_2 \rangle_{S_2}+\left(\frac{n_{2,1}}{p^{S_2+1}}+\frac{p-1}{p^{S_2+2}}+ \cdots + \frac{p-1}{p^{S_2+e}},\frac{n_{2,2}}{p^{S_2+1}}, \ldots, \frac{n_{2,l_2}}{p^{S_2+1}} \right)\\
         \vdots \\
		v_r=\langle \rho_r \rangle_{S_r}+\left(\frac{n_{r,1}}{p^{S_r+1}}+\frac{p-1}{p^{S_r+2}}+ \cdots + \frac{p-1}{p^{S_r+e}},\frac{n_{r,2}}{p^{S_r+1}}, \ldots, \frac{n_{r,l_r}}{p^{S_r+1}} \right).
\end{cases}
\] 
In this way $v =(v_1,\ldots,v_r) \in \RR_{\geq 0}^{l_1+\cdots+l_r}$ verifies $p^{S+e}v \in \ZZ_{\geq 0}^{l_1+\cdots+l_r}$ and moreover
\begin{align*}
|v|&=\sum_{i=1}^{r}\left(|\langle \rho_i \rangle_{S_i}|+\frac{p^{e}-1}{p^{S_i+e}} \right).
\end{align*}
Observe also that  for $i\in\{1,\ldots,r\}$, $v_i \preceq \langle \rho_i \rangle_{S_i+e} \preceq \langle \rho_i \rangle_{S+e}$ and by Proposition~\ref{propoBinomio}(2) 
\[\prod_{i=1}^{r}\binom{p^{S+e}|v_i|}{p^{S+e}v_i}c_{a(i)}^{d_{e}(i)} \cdot  \prod_{i=r+1}^{t}\binom{p^{S+e}|\langle \rho_i \rangle_{S+e}|}{p^{S+e}\langle \rho_i \rangle_{S+e}}c_{a(i)}^{b_{e}(i)}
\]
appears as the coefficient of the monomial 
\[(x^{p^{S+e}  \widehat{\E}_1 v_1}  \cdots  x^{p^{S+e}\widehat{\E}_r v_r})\cdot (  x^{p^{S+e}\widehat{\E}_{r+1} \langle \rho_{r+1} \rangle_{S+e}}\cdots   x^{p^{S+e}\widehat{\E}_t \langle \rho_t \rangle_{S+e} })
\]
in the polynomial $g$ given by the product 
$\prod_{i=1}^{r} f_i^{p^{S+e}|v_i|} \cdot \prod_{i =r+1}^t  f_i^{p^{S+e}|\langle \rho_i \rangle_{S+e}|}$. By the definitions of  $S$ and $S_i$, we conclude from Lucas's Lemma that $\binom{p^{S+e}|v_i|}{p^{S+e}v_i} \not \equiv 0$ mod $p$ for each $v_i$. Since $S_i=\infty$, 
for $ i \in \{r+1,\ldots, t\}$, then $\binom{p^{S+e}|\langle \rho_i \rangle_{S+e}|}{p^{S+e}\langle \rho_i \rangle_{S+e}} \not \equiv 0$ mod $p$. In conclusion one can set $Q(e)=p^{S+e} |v| +p^{S+e}|\langle( \rho_{r+1},\ldots, \rho_t) \rangle_{S+e}|$ to verify that 
$g \in  \fa^{Q(e)}$ but $g \notin \m^{[p^{S+e}]}$. In this way $Q(e)\leq \nu_{\fa}(p^{S+e})$. Dividing by $p^{S+e}$ and letting $e \rightarrow \infty$ yields finally  
\[\sum_{i=1}^{r}\left(|\langle \rho_i \rangle_{S_i}|+  \frac{1}{p^{S_i}}\right) + \sum_{i=r+1}^t|\rho_i| \leq \fpt(\fa).
\]
\end{proof}

\end{document}
